Table~\ref{tab:families} and Fig.~\ref{fig:campaign}(b) compare six families
run under the identical protocol, so that the only difference between two rows
is the space in which the seven algorithms search.

\paragraph{Same structure, different coordinates} PIDF-direct searches the
same 2-DOF PIDF ($b=1$, $c=0$, $N=N_C$) directly in
$(\log K_p,\log K_i,\log K_d,\log K_b)$ over the gain ranges spanned by the LQI
box. It reaches the same optimum (best $J_s=\BestPIDF$ against \CMObest{}) with
a similar number of feasible runs (\FeasPIDF{} against \CMOfeas{} of 180), and
the paired comparison on the same (algorithm, seed) cells is not significant
($p=\AbPIDFp$). The LQI coordinates therefore do not make the search easier
on this problem. What they change is what every candidate is. Inverting the map
\eqref{eq:map} and applying Kalman's return-difference condition shows that the
best direct-search design is itself LQ-optimal, $K=\IObestK$ with
$\min_\omega|1+L_K(j\omega)|=1$, i.e.\ the gain search converges onto the LQ
manifold; yet only \IOk{} of its \IOn{} feasible designs satisfy the condition,
so \IOpct\,\% of the controllers an equally lucky gain search may return carry
no LQ optimality, no guaranteed margins and no reason to be certifiable. The
LQI parameterization confines the search to that manifold from the first
evaluation, at no cost in optimality.

\paragraph{Richer and poorer structures} PID-direct (backward-difference
derivative) reaches best $J_s=\BestPID$ in \FeasPID{} feasible runs and is not
significantly different from CMO--LQI--PIDF ($p=\AbPIDp$). PIDF7-direct, with
free filter pole and setpoint weights, finds more feasible runs (\FeasPIDFs)
but a worse optimum (\BestPIDFs), with $N=\PIDFsN$\,rad/s, $b=\PIDFsb$ and
$c=\PIDFsc$; on the switched mission this design holds the \SI{1}{V} level
with \SwPIDFsEss{}\,mV of error and IAE \SwPIDFsIAE{}\,V\,s, i.e.\ the extra
freedom buys switched robustness at the price of the averaged objective. PI-direct never becomes feasible (minimum
violation \ViolPI{}): a PI cannot meet the overshoot, load-deviation and
current specifications of the six scenarios at once, and the paired comparison
is decisive ($p=\AbPIp$).

\paragraph{The full-state twin} CMO--LQI--SF uses the same four variables but
applies the gain as a measured-state LQI. It is feasible more often
(\FeasSF{} runs) but at a worse averaged optimum (\BestSF) than its
output-feedback realization ($p=\AbSFp$ in favour of the PIDF on the averaged
objective). On the switched mission the ordering reverses
(Section~\ref{sec:res:switched}): the twin is the more robust controller at the
\SI{1}{V} level, which is the switched counterpart of the output-current term
of \eqref{eq:gap}.

\paragraph{What the protocol, rather than the parameterization, buys}
Every family tuned under the constrained multi-fidelity protocol ends close to
full compliance on the reference mission (Table~\ref{tab:switched}: $u$ between
\SwPIdUtil{} and \SwUtil{}; the lowest value belongs to a slow PI whose IAE,
\SwPIdIAE{}\,V\,s, is 2.6 times that of the others), whereas the six fixed-rule designs lie between
\SwLQIUtil{} and \SwLQRUtil{}. The decisive ingredients are the
engineering-unit specification and the switched transition screen; the choice
of search coordinates decides the guarantees that come with the result.
